\makeatletter \def\input@path{{../../}} \makeatother
\documentclass[11pt]{article}

\usepackage[utf8]{inputenc}
\usepackage{graphicx}
\usepackage[dvipsnames]{xcolor}
\usepackage{color}
\usepackage[margin=1in]{geometry}
\usepackage{hyperref}
\usepackage{tikz}
\usepackage{amsmath}
\usepackage{commath}
\usepackage{relsize}
\usepackage{centernot}
\usepackage{amssymb}
\usepackage{amsfonts}
\usepackage{graphicx}
\usepackage{textcomp}
\usepackage{gensymb}
\usepackage{enumitem}
\usepackage{siunitx}
\usepackage{fancyhdr}
\usepackage{floatrow}
\usepackage{wrapfig}
\usepackage[labelfont=bf]{caption}
\usepackage{mathtools}
\usepackage{pgfplots}
\usepackage{multicol}
\usepackage[misc]{ifsym}
\usepackage{tabularx}
\usepackage{empheq}
\usepackage{tkz-tab}
\usepackage{xpatch}
\usepackage{caption}
\usepackage{subcaption}
\usepackage{tabularray}
\usepackage{esint}
\usepackage{amsthm}
\usepackage{thmtools}
\RequirePackage[most]{tcolorbox}

\swapnumbers

\makeatletter
\declaretheoremstyle[
  headfont= \sffamily\bfseries\color{cyan!50!black},
  headpunct={\sffamily\bfseries.},
  postheadspace=0.5em,
  notefont=\sffamily\bfseries,
  headformat=\NAME\NUMBER\let\thmt@space\@empty\NOTE,
  bodyfont=\mdseries\slshape,
  spaceabove=12pt,spacebelow=12pt,
]{thmstyle}
\makeatother

\theoremstyle{thmstyle}
\declaretheorem[name=Definición]{theorem}
\newtheorem{definition}[theorem]{Definición}

\definecolor{lightcyan}{rgb}{0.88, 1.0, 1.0}
\colorlet{mythmback}{lightcyan!40!white}

\tcolorboxenvironment{theorem}{
  colback=mythmback,coltitle=blue,colframe=mythmback,
  left=0pt,right=0pt,
  top=0pt,bottom=0pt,
  before skip=10pt, after skip=10pt,
}

\xpatchcmd{\tkzTabLine}{$0$}{$\bullet$}{}{}
\tikzset{t style/.style={style=solid}}
\pgfplotsset{compat=newest}

\newtheorem{manualtheoreminner}{}
\newenvironment{manualtheorem}[1]{%
  \renewcommand\themanualtheoreminner{#1}%
  \manualtheoreminner
}{\endmanualtheoreminner}

\renewcommand{\pd}[2]{\frac{\partial #1}{ \partial #2}}
\newcommand{\td}[2]{\frac{\mathrm{d} #1}{ \mathrm{d} #2}}
\newcommand{\pdd}[2]{\frac{\partial^2 #1}{ \partial #2 ^2}}
\newcommand{\pddc}[3]{\frac{\partial^2 #1}{ \partial #2 \partial #3}}
\newcommand{\solution}[0]{\large{\textbf{Solución}}\normalsize}

\def\sca   #1{\mbox{\rm{#1}}{}}
\def\mat   #1{\mbox{\boldmath $\mathsf #1$}}
\def\vec   #1{\mbox{\boldmath $#1$}{}}
\def\ten   #1{\mbox{\boldmath $#1$}{}}
\def\ltr   #1{\mbox{\sf{#1}}}
\def\bltr  #1{\mbox{\sffamily{\bfseries{{#1}}}}}

\DeclareMathOperator{\dive}{div}
\DeclareMathOperator{\trace}{trace}
\DeclareMathOperator{\tr}{tr}
\DeclareMathOperator{\symm}{symm}
\DeclareMathOperator{\sk}{skew}
\DeclareMathOperator{\grad}{grad}
\DeclareMathOperator{\Grad}{Grad}
\DeclareMathOperator{\curl}{curl}
\DeclareMathOperator{\Curl}{Curl}
\DeclareMathOperator*{\argmin}{\arg\!\min}

\def\dx{\mbox{\(\,\mathrm{d}x\)}}
\def\ds{\mbox{\(\,\mathrm{d}s\)}}
\def\dS{\mbox{\(\,\mathrm{d}S\)}}
\def\dV{\mbox{\(\,\mathrm{d}V\)}}
\def\dv{\mbox{\(\,\mathrm{d}v\)}}
\def\R{\mathbb R}
\def\C{\mathbb C}
\def\N{\mathbb N}
\def\Q{\mathbb Q}
\def\Z{\mathbb Z}
\def\D{\mathcal D}

\usepackage{mdframed}
\mdfdefinestyle{MyFrame}{%
    linecolor=black,
    linewidth=1pt,
    outerlinewidth=1pt,
    roundcorner=20pt,
    innertopmargin=\baselineskip,
    innerbottommargin=\baselineskip,
    innerrightmargin=20pt,
    innerleftmargin=20pt,
    splittopskip=2\baselineskip,
    backgroundcolor=gray!50!white}

\setlength\textwidth{7in}
\setlength\parindent{0pt}
\oddsidemargin=-0.5cm
\pagestyle{fancy}
\newcommand*{\vertbar}{\rule[-1ex]{0.5pt}{2.5ex}}
\setlength{\headheight}{13.59999pt}

\renewcommand{\headrulewidth}{0.1pt}
\renewcommand{\footrulewidth}{0.1pt}
\renewcommand{\figurename}{Figura}
\renewcommand\refname{Referencias}
\renewcommand{\thesection}{\arabic{section}.}
\renewcommand{\thesubsection}{\thesection\arabic{subsection}.}
\renewcommand{\labelitemi}{{\raisebox{.3\height}{\scalebox{0.6}{$\blacklozenge$}}}}
\renewcommand*\contentsname{\'Indice}

\newcommand{\p}{\vspace{10pt}}

\AtBeginDocument{\vspace*{-3.2\baselineskip}}
\textheight=665pt

\lhead{\scshape Pontificia Universidad Cat\'olica de Chile}
\rhead{\scshape IMC UC}
\cfoot{\thepage}

\usepackage{footnote}
\usepackage{etoolbox}
\newtoggle{indefintion}
\togglefalse{indefintion}
\pretocmd{\footnote}{\iftoggle{indefintion}{\stepcounter{footnote}}{\relax}}{}{}

\begin{document}
    \begin{flushleft}
        \small
        \vspace{1pt}
        \textbf{IMT3130} \hfill
        \textbf{10/04/2026}
    \end{flushleft}
    \begin{center}
        \LARGE\textbf{Ayudant\'ia 4}\\
        \vspace{3pt}
        \normalsize\textbf{Formulaciones d\'ebiles, Poincar\'e, Lax--Milgram y Galerkin}\\
        \normalsize Ayudante: Basti\'an Herrera \Letter{\hspace{1pt}: \texttt{biherrera@uc.cl}}
        \rule{\textwidth}{0.05mm}
    \end{center}

\section*{Problemas}

\begin{enumerate}[label=\textbf{\arabic*.}, leftmargin=*]

    \item \textbf{Problema con condiciones mixtas.} Sea $\Omega\subset\R^d$ un dominio acotado, conexo y Lipschitz. Suponga que $\partial\Omega=\overline{\Gamma_D}\cup\overline{\Gamma_R}$, con $\Gamma_D\cap\Gamma_R=\varnothing$. Para $c\in L^\infty(\Omega)$, $\beta\in L^\infty(\Gamma_R)$ y $c,\beta\ge 0$, considere el problema
    \[
    \begin{aligned}
        -\Delta u + cu &= f &&\text{en }\Omega,\\
        \gamma_0u &=0 &&\text{en }\Gamma_D,\\
        \partial_n u+\beta\gamma_0u &= r &&\text{en }\Gamma_R.
    \end{aligned}
    \]
    Defina
    \[
        V=\{v\in H^1(\Omega):\gamma_0v=0\text{ en }\Gamma_D\}.
    \]
    \begin{enumerate}[label=(\alph*)]
        \item Derive la formulaci\'on d\'ebil del problema, identificando la forma bilineal $a:V\times V\to\R$ y la forma lineal $\ell:V\to\R$. Explique por qu\'e la condici\'on de Dirichlet es esencial y la de Robin es natural.
        \item Bajo hip\'otesis razonables sobre $f$ y $r$, demuestre continuidad de $a$ y de $\ell$ en $V$. Luego determine en qu\'e casos $a$ es coerciva en $V$ y concluya existencia y unicidad por Lax--Milgram.
        \item Considere ahora al caso de condiciones puramente Neumann, i.e. $\Gamma_D=\varnothing$, $\Gamma_R=\partial\Omega$, $\beta=0$ y $c=0$. Muestre que las funciones constantes pertenecen a $\ker a$, deduzca la condici\'on necesaria de compatibilidad tomando $v=1$, y explique c\'omo se recupera coercividad sobre
        \[
            H^1_\ast(\Omega)=\left\{v\in H^1(\Omega):\int_\Omega v\,dx=0\right\}
        \]
        usando la desigualdad de Poincar\'e--Wirtinger. ¿Qu\'e cambia si $c(x)\ge c_0>0$?
    \end{enumerate}

    \item \textbf{Ortogonalidad de Galerkin como proyecci\'on.} Sea $V$ un espacio de Hilbert, $a:V\times V\to\R$ una forma bilineal sim\'etrica, continua y coerciva, y $\ell\in V'$. Sea $u\in V$ la soluci\'on de
    \[
        a(u,v)=\ell(v)
        \qquad\forall v\in V.
    \]
    Sea $V_h\subset V$ un subespacio finito-dimensional conforme y suponga que $u_h\in V_h$ satisface
    \[
        a(u_h,v_h)=\ell(v_h)
        \qquad\forall v_h\in V_h.
    \]
    \begin{enumerate}[label=(\alph*)]
        \item Demuestre la ortogonalidad de Galerkin
        \[
            a(u-u_h,v_h)=0
            \qquad\forall v_h\in V_h.
        \]
        \item Defina $\|v\|_a\coloneqq a(v,v)^{1/2}$ y pruebe que es una norma equivalente a la norma de $V$.
        \item Use la ortogonalidad para concluir que $u_h$ es la mejor aproximaci\'on de $u$ en $V_h$ respecto de $\|\cdot\|_a$, es decir,
        \[
            \|u-u_h\|_a=\inf_{v_h\in V_h}\|u-v_h\|_a.
        \]
        \item Si $V_h=\operatorname{span}\{\varphi_1,\dots,\varphi_N\}$ y $u_h=\sum_{j=1}^N U_j\varphi_j$, escriba el sistema lineal que satisfacen los coeficientes $U_j$.
    \end{enumerate}

\end{enumerate}

\end{document}
